\documentclass[a4paper,12pt]{article}
\usepackage[utf8]{inputenc}
\usepackage[T1]{fontenc}
\usepackage[ngerman]{babel}
\usepackage{amsmath, amssymb}
\usepackage{geometry}
\usepackage{tikz}
\usepackage{tikz-cd}
\usetikzlibrary{shapes.geometric, arrows.meta, positioning, fit, backgrounds}
\usepackage{parskip}
\usepackage{titlesec}
\usepackage{booktabs}
\usepackage{enumitem}

\geometry{left=3cm, right=3cm, top=2.5cm, bottom=2.5cm}

\titleformat{\section}{\large\bfseries}{}{0em}{}[\titlerule]
\titleformat{\subsection}{\normalsize\bfseries}{}{0em}{}

\tikzset{
  box/.style={rectangle, draw, rounded corners=3pt, text width=4.5cm,
              align=center, minimum height=1.0cm, font=\small, inner sep=4pt},
  boxw/.style={box, text width=5.5cm},
  arrow/.style={-{Latex[length=2mm]}, thick},
}

\begin{document}

\begin{center}
  {\LARGE\bfseries Studierhinweise zu Lektion 7}\\[0.5em]
  {\large Mathematische Grundlagen (Modul 61111)}
\end{center}

\vspace{1em}

Diese Lektion ist zweigeteilt. Das erste Kapitel schließt den Analysis-Teil des
Moduls ab: Es wird das \textbf{Riemann-Integral} eingeführt, das die intuitive
Vorstellung von Flächeninhalten präzisiert. Anschließend verbindet der erste und
zweite \textbf{Hauptsatz der Differential- und Integralrechnung} das Integrieren
mit dem Differenzieren, was die praktische Berechnung von Integralen erst
ermöglicht.

Der weitaus umfangreichere zweite Teil führt in die \textbf{formale Logik} ein.
Die Aussagenlogik und die Prädikatenlogik sind künstliche Sprachen, die
mathematische Sachverhalte ohne Mehrdeutigkeiten ausdrücken. Sie bilden die
Grundlage für Anwendungen in der Digitaltechnik, Komplexitätstheorie,
Datenbanken und Logikprogrammierung.

Lesen Sie mathematische Texte nie passiv: Vollziehen Sie jeden Schritt nach,
konstruieren Sie eigene Beispiele, und prüfen Sie, ob Sie Definitionen aus dem
Gedächtnis reproduzieren können -- gerade bei der Logik, wo die Präzision der
Formulierungen entscheidend ist.

% -----------------------------------------------------------------------
\section*{Struktur der Lektion 7}

\begin{center}
\begin{tikzpicture}[node distance=0.6cm and 1.0cm]

  % Column headers
  \node[font=\bfseries\small] (h1) at (0,0)   {Kapitel 20: Integration};
  \node[font=\bfseries\small] (h2) at (8.0,0) {Kapitel 21: Logik};

  % Kapitel 20 boxes
  \node[box, below=0.8cm of h1] (ri)  {20.1 Riemann-Integral\\Partition, Unter-/Obersumme,\\Integrierbarkeit};
  \node[box, below=0.5cm of ri]  (hs)  {20.2 Hauptsätze\\Stammfunktion, partielle\\Integration, Substitution};

  % Kapitel 21: two columns
  \node[box, text width=3.8cm, below=0.8cm of h2, xshift=-2.3cm] (al1) {21.1.1 Syntax\\Atome, Formeln,\\Präzedenz\-regeln};
  \node[box, text width=3.8cm, below=0.8cm of h2, xshift=+2.3cm] (pl1) {21.2.2 Syntax\\Terme, Formeln,\\Quantoren};
  \node[box, text width=3.8cm, below=0.5cm of al1] (al2) {21.1.2 Semantik\\Bewertung, Wahr\-heitstafel,\\Tautologie, Folgerung};
  \node[box, text width=3.8cm, below=0.5cm of pl1] (pl2) {21.2.3 Semantik\\Interpretation,\\Universum, Wahrheit};
  \node[box, text width=3.8cm, below=0.5cm of al2] (al3) {21.1.3 Normalformen\\NNF, KNF, DNF};
  \node[box, text width=3.8cm, below=0.5cm of pl2] (pl3) {21.2.4 Normalformen\\NNF, pränexe NF};
  \node[box, text width=3.8cm, below=0.5cm of al3] (al4) {21.1.4 Beweise\\Äquivalenz- und\\Schlussregeln};
  \node[box, text width=3.8cm, below=0.5cm of pl3] (pl4) {21.2.5 Beweise\\$\forall$-i, $\exists$-i,\\Syllogismus};

  % Kapitel 21 header labels
  \node[font=\small\bfseries, above=0.3cm of al1] {21.1 Aussagenlogik};
  \node[font=\small\bfseries, above=0.3cm of pl1] {21.2 Prädikatenlogik};

  % Arrows Kap 20
  \draw[arrow] (ri) -- (hs);

  % Arrows AL
  \draw[arrow] (al1) -- (al2);
  \draw[arrow] (al2) -- (al3);
  \draw[arrow] (al3) -- (al4);

  % Arrows PL
  \draw[arrow] (pl1) -- (pl2);
  \draw[arrow] (pl2) -- (pl3);
  \draw[arrow] (pl3) -- (pl4);

  % Cross arrow: AL → PL
  \draw[arrow] (al1.east) -- ++(0.5,0) |- (pl1.west);

  % Cross arrow: Kap20 → Kap21
  \draw[arrow] (hs.east) -- ++(0.4,0) |- (al1.west);

\end{tikzpicture}
\end{center}

\newpage

% -----------------------------------------------------------------------
\section*{Zielelement 20.1 -- Das Riemann-Integral}

\subsection*{Lerninhalte}

\begin{center}
\begin{tikzpicture}[node distance=0.6cm and 1.0cm]
  \node[box] (par)  {Partition $P$ von $[a,b]$\\$a = t_0 < t_1 < \cdots < t_n = b$};
  \node[box, below=0.5cm of par] (us)  {Untersumme $U(f,P)$\\Obersumme $O(f,P)$};
  \node[box, below=0.5cm of us]  (vf)  {Verfeinerung $Q$ von $P$:\\$U(f,P) \le U(f,Q)$,\; $O(f,P) \ge O(f,Q)$};
  \node[box, below=0.5cm of vf]  (int) {Integrierbarkeit:\\$\sup U = \inf O$\\Riemann-Integral $\int_a^b f$};
  \node[box, below=0.5cm of int] (rec) {Rechenregeln:\\Linearität, Additivität,\\Abschätzung $m(b-a) \le \int \le M(b-a)$};
  \node[box, below=0.5cm of rec] (ubs) {Unbestimmtes Integral\\$F(x) = \int_a^x f(t)\,dt$ ist stetig};
  \draw[arrow] (par) -- (us);
  \draw[arrow] (us)  -- (vf);
  \draw[arrow] (vf)  -- (int);
  \draw[arrow] (int) -- (rec);
  \draw[arrow] (rec) -- (ubs);
\end{tikzpicture}
\end{center}

\subsection*{Lernziele}

Nach Durcharbeiten dieses Abschnitts sollten Sie:
\begin{itemize}
  \item den Begriff der Partition sowie von Unter- und Obersumme erklären und
        für konkrete Funktionen berechnen können,
  \item das Verhalten von Unter-/Obersummen bei Verfeinerungen verstehen
        (Lemma~20.1.5) und begründen können,
  \item die Definition der Riemann-Integrierbarkeit (Definition~20.1.7) kennen
        und die $\varepsilon$-Charakterisierung (Proposition~20.1.9) anwenden können,
  \item die Rechenregeln der Integration (Linearität, Additivität über Intervalle,
        Abschätzung) sicher einsetzen können,
  \item die Stetigkeit des unbestimmten Integrals $F(x) = \int_a^x f(t)\,dt$
        begründen können (Satz~20.1.14).
\end{itemize}

\subsection*{Selbstkontrollelement 20.1}

Zeigen Sie, dass die Funktion $f : [0,1] \to \mathbb{R}$ mit
\[
  f(x) = \begin{cases} 1, & x \in \{0, \tfrac{1}{2}, 1\}, \\ 0, & \text{sonst} \end{cases}
\]
nach Proposition~20.1.9 integrierbar ist, und bestimmen Sie $\int_0^1 f(x)\,dx$.

\newpage

% -----------------------------------------------------------------------
\section*{Zielelement 20.2 -- Hauptsätze der Differential- und Integralrechnung}

\subsection*{Lerninhalte}

\begin{center}
\begin{tikzpicture}[node distance=0.6cm and 1.0cm]
  \node[box] (hs1) {1.\ Hauptsatz:\\$F(x)=\!\int_a^x\! f(t)\,dt \Rightarrow F'(c)=f(c)$\\(falls $f$ in $c$ stetig)};
  \node[box, below=0.5cm of hs1] (stf) {Stammfunktion $g$ mit $g'=f$\\Korollar: stetig $\Rightarrow$ integrierbar};
  \node[box, below=0.5cm of stf] (hs2) {2.\ Hauptsatz:\\$\int_a^b f = g(b) - g(a)$};
  \node[box, below=0.5cm of hs2] (tab) {Tabelle der Stammfunktionen\\(Polynome, $\exp$, $\ln$, $\sin$, $\cos$, \ldots)};
  \node[box, below=0.5cm of tab] (pi)  {Partielle Integration\\$\int f g' = fg\big|_a^b - \int f' g$};
  \node[box, below=0.5cm of pi]  (sub) {Substitutionsregel\\$\int_a^b f(g(x))g'(x)\,dx = \int_{g(a)}^{g(b)} f(u)\,du$};
  \draw[arrow] (hs1) -- (stf);
  \draw[arrow] (stf) -- (hs2);
  \draw[arrow] (hs2) -- (tab);
  \draw[arrow] (tab) -- (pi);
  \draw[arrow] (pi)  -- (sub);
\end{tikzpicture}
\end{center}

\subsection*{Lernziele}

Nach Durcharbeiten dieses Abschnitts sollten Sie:
\begin{itemize}
  \item den ersten und zweiten Hauptsatz der Differential- und Integralrechnung
        formulieren und beweisen können,
  \item aus dem ersten Hauptsatz folgern können, dass jede stetige Funktion eine
        Stammfunktion besitzt (Korollar~20.2.4),
  \item den zweiten Hauptsatz zur Berechnung von Integralen anwenden können,
        indem Sie eine Stammfunktion bestimmen,
  \item die Stammfunktionen elementarer Funktionen aus der Tabelle~20.2.7
        abrufen und einsetzen können,
  \item partielle Integration und Substitutionsregel korrekt anwenden können,
        insbesondere das richtige Identifizieren von $f$, $g$ bzw.\ $g'$.
\end{itemize}

\subsection*{Selbstkontrollelement 20.2}

Berechnen Sie die folgenden Integrale:
\begin{enumerate}[label=(\alph*)]
  \item $\displaystyle\int_0^1 x e^x\,dx$ \quad (partielle Integration)
  \item $\displaystyle\int_0^{\pi/2} \sin^3(x)\cos(x)\,dx$ \quad (Substitution)
\end{enumerate}

\newpage

% -----------------------------------------------------------------------
\section*{Zielelement 21.1 -- Aussagenlogik}

\subsection*{Lerninhalte}

\begin{center}
\begin{tikzpicture}[node distance=0.6cm and 1.0cm]
  \node[box] (syn)  {Syntax:\\Atome, Formeln induktiv,\\Junktoren $\lnot,\land,\lor,\to,\leftrightarrow$};
  \node[box, below=0.5cm of syn] (pra) {Präzedenz- und\\Klammerregeln};
  \node[box, below=0.5cm of pra] (sem) {Semantik: Bewertung $I$\\Wahrheitstafel, Tautologie,\\Widerspruch, Erfüllbarkeit};
  \node[box, below=0.5cm of sem] (fol) {Semantische Folgerung $\alpha \models \beta$\\Logische Äquivalenz $\alpha \approx \beta$};
  \node[box, below=0.5cm of fol] (nor) {Normalformen:\\NNF, KNF, DNF};
  \node[box, below=0.5cm of nor] (bew) {Formale Beweise:\\Äquivalenz- und Schlussregeln\\(mp, mt, con, simp, add)};
  \draw[arrow] (syn) -- (pra);
  \draw[arrow] (pra) -- (sem);
  \draw[arrow] (sem) -- (fol);
  \draw[arrow] (fol) -- (nor);
  \draw[arrow] (nor) -- (bew);
\end{tikzpicture}
\end{center}

\subsection*{Lernziele}

Nach Durcharbeiten dieses Abschnitts sollten Sie:
\begin{itemize}
  \item den induktiven Aufbau aussagenlogischer Formeln (Definition~21.1.1)
        erklären und konkrete Ausdrücke auf ihre Korrektheit prüfen können,
  \item Präzedenzregeln anwenden und überflüssige Klammern entfernen können,
  \item einer Formel für eine gegebene Bewertung einen Wahrheitswert zuordnen
        und vollständige Wahrheitstafeln aufstellen können,
  \item die Begriffe tautologisch, widerspruchsvoll, erfüllbar und falsifizierbar
        unterscheiden und anhand von Wahrheitstafeln nachweisen können,
  \item semantische Folgerung ($\alpha \models \beta$) und logische Äquivalenz
        ($\alpha \approx \beta$) definieren und die Zusammenhänge aus Bemerkung~21.1.22
        kennen,
  \item die De-Morganschen Gesetze und weitere Äquivalenzregeln (21.1.26) anwenden,
  \item Formeln in Negationsnormalform, konjunktive und disjunktive Normalform
        überführen können,
  \item formale Beweisfolgen mit den Äquivalenz- und Schlussregeln aufstellen können,
        insbesondere Modus ponens und Modus tollens sicher einsetzen.
\end{itemize}

\subsection*{Selbstkontrollelement 21.1}

Sei $\alpha = (A \to B) \land (B \to C)$.
\begin{enumerate}[label=(\alph*)]
  \item Stellen Sie die vollständige Wahrheitstafel von $\alpha$ auf.
  \item Ist $\alpha$ eine Tautologie? Begründen Sie.
  \item Leiten Sie mit einer formalen Beweisfolge aus $\alpha$ und der Prämisse $A$
        die Formel $C$ her.
  \item Überführen Sie $\alpha$ in KNF.
\end{enumerate}

\newpage

% -----------------------------------------------------------------------
\section*{Zielelement 21.2 -- Prädikatenlogik}

\subsection*{Lerninhalte}

\begin{center}
\begin{tikzpicture}[node distance=0.6cm and 1.0cm]
  \node[box] (trm)  {Terme: Variablen $V$,\\Konstanten $K$, Funktions\-symbole $F$};
  \node[box, below=0.5cm of trm] (frm) {Formeln: Prädikate $P(t_1,\ldots,t_n)$,\\Gleichung $t_1=t_2$, Junktoren,\\Quantoren $\forall$, $\exists$};
  \node[box, below=0.5cm of frm] (bnd) {Freie / gebundene Variablen\\Wirkungsbereich eines Quantors\\Konsistente Umbenennung};
  \node[box, below=0.5cm of bnd] (int) {Interpretation $(U, I)$\\Universum, Zuordnung\\zu Konstanten, Funktionen, Prädikaten};
  \node[box, below=0.5cm of int] (nor) {Normalformen:\\NNF, pränexe Normalform (PNF)};
  \node[box, below=0.5cm of nor] (bew) {Formale Beweise:\\$\forall$-Instantiierung, $\exists$-Instantiierung\\Syllogismus (Aristoteles)};
  \draw[arrow] (trm) -- (frm);
  \draw[arrow] (frm) -- (bnd);
  \draw[arrow] (bnd) -- (int);
  \draw[arrow] (int) -- (nor);
  \draw[arrow] (nor) -- (bew);
\end{tikzpicture}
\end{center}

\subsection*{Lernziele}

Nach Durcharbeiten dieses Abschnitts sollten Sie:
\begin{itemize}
  \item Terme und prädikatenlogische Formeln induktiv aufbauen und erkennen können,
  \item freie und gebundene Vorkommen von Variablen in einer Formel bestimmen
        und den Wirkungsbereich jedes Quantors angeben können,
  \item eine Formel konsistent umbenennen können (Merkregel~21.2.17),
  \item einer Formel für eine konkrete Interpretation einen Wahrheitswert zuordnen
        (Definition~21.2.22),
  \item die Äquivalenzregeln für Quantoren (Proposition~21.2.32) kennen und
        anwenden können: Quantorwechsel, Quantortausch, Quantorzusammenfassung,
        Quantorelimination, Quantifizierung,
  \item Formeln in die Negationsnormalform und die pränexe Normalform überführen,
  \item einfache prädikatenlogische Argumente mit $\forall$-Instantiierung und
        $\exists$-Instantiierung formal beweisen können.
\end{itemize}

\subsection*{Selbstkontrollelement 21.2}

Sei $\alpha = \forall x \exists y\, P(x, y)$.
\begin{enumerate}[label=(\alph*)]
  \item Bestimmen Sie für jede der folgenden Interpretationen den Wahrheitswert
        von $\alpha$:
        \begin{enumerate}[label=(\roman*)]
          \item $U = \mathbb{N}$,\quad $I(P) = \{(m,n) \mid m < n\}$.
          \item $U = \{0\}$,\quad $I(P) = \emptyset$.
        \end{enumerate}
  \item Begründen Sie, warum $\forall x \exists y\, P(x,y)$ im Allgemeinen
        \emph{nicht} äquivalent zu $\exists y \forall x\, P(x,y)$ ist.
  \item Überführen Sie die Formel $\neg(\forall x\, P(x) \land \exists y\, Q(y))$
        in Negationsnormalform.
\end{enumerate}

\end{document}
